\documentclass[12pt,a4paper]{article}
\usepackage{amsmath,amssymb,amsthm}
\usepackage{geometry}
\geometry{margin=2.5cm}
\usepackage{hyperref}

\title{\bf Steps to Prove the Spectral Determinant Identity\\ for the Generalised Helix}
\author{}
\date{}

\begin{document}
\maketitle

Let \(L(s)\) be an automorphic \(L\)-function with completed form \(\Lambda(s)\), functional equation, and explicit formula
\[
\sum_{\rho} h(\rho) = \sum_{p^m} \frac{c(p^m)}{p^{m/2}}\,h(m\log p) + \text{archimedean terms},
\]
where \(c(p^m)\) are the arithmetic coefficients.  The operator \(D_{L,\lambda}\) is the Dirac operator on the circle \(\mathbb{T}_L\) (\(L = 2\log\lambda\)) with point interactions at all \(m\log p \bmod L\) (\(p^m \le \lambda^2\)) and coupling constants \(\alpha_{p^m} = c(p^m)/p^{m/2}\).  The limiting operator on \(\mathbb{R}\) is \(D_{L,\infty}\).

We outline the five steps to prove
\[
\lim_{\lambda\to\infty} \Delta_\lambda(z)\,\Gamma_{\text{arch}}(z) = \Lambda(z),
\]
where \(\Delta_\lambda(z) = \det_{\text{reg}}\bigl((D_{L,\lambda} - z)(D_0 - z)^{-1}\bigr)\) and \(\Gamma_{\text{arch}}\) is the archimedean correction.

\begin{enumerate}

\item \textbf{Explicit finite determinant.}  
Using the Krein resolvent formula for point interactions on a circle, write
\[
\Delta_\lambda(z) = \det\bigl(1 + K_\lambda(z)\bigr),
\]
where \(K_\lambda(z)\) is a finite‑rank operator encoding the coupling constants and the free Dirac propagator between the points \(x_j = m\log p \bmod L\).  The determinant is an entire function of \(z\) whose zeros are the eigenvalues of \(D_{L,\lambda}\).

\item \textbf{Relation to truncated zeros.}  
By the zero‑locking condition \(\phi_L(\gamma_n) = 2\pi n\), the eigenvalues of \(D_{L,\lambda}\) approximate the zeros \(\gamma_n\) of \(L(s)\) up to height \(O(\lambda^2)\).  Consequently,
\[
\Delta_\lambda(z) \;=\; e^{A(z)} \prod_{|\gamma_n| \le T(\lambda)} \Bigl(1 - \frac{z}{\gamma_n}\Bigr) e^{z/\gamma_n},
\]
where \(A(z)\) is an entire function coming from the regularisation and \(T(\lambda) \to \infty\) as \(\lambda\to\infty\).  The factors \(e^{z/\gamma_n}\) ensure convergence.

\item \textbf{Comparison with the CCM matrix.}  
The truncated Weil quadratic form \(Q_L\) of the Connes–Consani–Moscovici spectral triple for \(L(s)\) is a finite Hermitian matrix whose elements are built from the arithmetic coefficients \(c(p^m)/p^{m/2}\).  Show that
\[
\det(Q_L - zI) \;=\; \Delta_\lambda(z) \cdot (\text{explicit Gamma factor}) \cdot (1 + o(1)),
\]
where the error tends to zero as \(\lambda \to \infty\).  This equivalence is established by comparing the boundary value problem of \(D_{L,\lambda}\) with the eigenvalue equation \(Q_L\xi = \epsilon\xi\) in the Fourier basis.

\item \textbf{Asymptotic identification with \(\Lambda(z)\).}  
Assume the generalised CCM Conjecture~1: the ground‑state eigenvector coefficients of \(Q_L\) are given by a prime sum with weights \(c(p^m)/p^{m/2}\) and a power‑saving error.  Then as \(\lambda\to\infty\),
\[
\Delta_\lambda(z)\,\Gamma_{\text{arch}}(z) \;\longrightarrow\; \Lambda(z)
\]
uniformly on compact sets, where \(\Gamma_{\text{arch}}(z)\) is the product of Gamma factors from the free Dirac operator on \(\mathbb{T}_L\) and the archimedean part of the functional equation of \(L(s)\).  The convergence follows from matching the Hadamard product of \(\Lambda(z)\) over zeros with the infinite‑product limit of the truncated zero products.

\item \textbf{Contour integral and trace formula.}  
For any test function \(h\) in the Weil class,
\begin{align*}
\operatorname{Tr}\bigl(h(D_{L,\infty})\bigr) - \operatorname{Tr}\bigl(h(D_0)\bigr)
&= \frac{1}{2\pi i} \oint_{\mathcal{C}} h(z) \frac{d}{dz}
\log\frac{\det_{\text{reg}}(D_{L,\infty} - z)}{\det_{\text{reg}}(D_0 - z)}\,dz \\
&= \frac{1}{2\pi i} \oint_{\mathcal{C}} h(z) \frac{\Lambda'(z)}{\Lambda(z)}\,dz + \text{trivial zeros} \\
&= \sum_{\rho} h(\rho) + \text{archimedean terms}
= \sum_{p^m} \frac{c(p^m)}{p^{m/2}} h(m\log p).
\end{align*}
Hence the trace formula holds exactly, and the spectrum of \(D_{L,\infty}\) coincides with the zeros of \(L(s)\).

\end{enumerate}

The main open difficulty lies in Step~4, which requires the generalised CCM Conjecture~1.  All other steps are technically involved but rest on well‑established operator theory and the geometry of the generalised helix.

\end{document}